\section{Construction principles and further questions}

The nineteen bounds arise from several distinct ways of enlarging a
constrained configuration. Tight contacts do not always prevent motion:
the 552-point layer in dimension 25 moves with its internal Gram matrix
unchanged. A saturated selection need not exhaust its candidate geometry:
recolouring changes the second-layer selection in dimension 27, while
shared blockers make a family of support insertions profitable. A local
capacity bound can be attained by importing a different combinatorial
object, as the weight-four designs and their Hadamard images show.

Two further mechanisms use information supplied by the ambient lattice.
Its minimum makes cross-shell compatibility automatic in dimension 38.
Its design moments determine a section count in dimension 43 and give
a witness inequality in dimension 45. In the latter case, the right
statistic changes the scale of the search: 298 explicit vectors certify
a configuration with more than seven million points. For dimensions
49--55, a common class supplies seven lifts, and exact union counts
distinguish the gain from class growth from the gain from image selection.

These results suggest structural questions as well as larger finite
searches. Which contacting layers admit a boundary-compatible motion?
When can a signed code attaining its local capacity be embedded in a
larger configuration? Which moment functionals isolate a count that is
accessible through a small witness set? Each question retains the
geometric constraints responsible for the results while opening choices
beyond the specific coordinates presented here.

Qiushi Engine developed these constructions through sustained autonomous
mathematical research, connecting geometric freedom, discrete selection
and exact counting. The resulting proofs and finite objects make those
connections available for further mathematics: they can be reconstructed,
varied and used as starting points for new constructions.
